\documentclass[11pt]{amsart}

\usepackage{amsmath}
\usepackage{amssymb}

%% Typographic only: relax line breaking slightly to avoid overfull
%% lines around long compound words and bracketed citations.
\emergencystretch=1.5em

%% ------------------------------------------------------------------
%% Theorem environments
%% ------------------------------------------------------------------
%% Results are numbered by section; where the source notes numbered a
%% result by its subsection (Lemma 4.3, Remark 4.5), the counter is
%% set explicitly so the printed numbers match the text's references.
\theoremstyle{plain}
\newtheorem{theorem}{Theorem}[section]
\newtheorem{lemma}{Lemma}[section]
\newtheorem{proposition}{Proposition}[section]
\newtheorem{corollary}{Corollary}[section]

\theoremstyle{remark}
\newtheorem{remark}{Remark}[section]

%% A proof environment whose head may wrap: amsthm's \proof puts the
%% optional title into an unbreakable \item label, which overflows the
%% line for the long proof head of Proposition 3.1. Same mechanics as
%% amsthm's proof, with the head set as ordinary (breakable) text.
\makeatletter
\newenvironment{longproof}[1]{\par
  \pushQED{\qed}%
  \normalfont \topsep6\p@\@plus6\p@\relax
  \trivlist
  \item\relax{\itshape #1\@addpunct{.}}\ \ignorespaces
}{%
  \popQED\endtrivlist\@endpefalse
}
\makeatother

%% ------------------------------------------------------------------

\begin{document}

\title[Certificates for pure cubic Thue equations]{Certificates for
  pure cubic Thue equations: Merlin--Arthur membership and the Compact
  Vanishing Problem}

\author{The diophantine-fable expedition}
\thanks{Draft --- authorship placeholder.}

\date{Sequel draft of 2026-07-22, assembled from the certificate
  notes; a companion to the exponential-time paper \cite{companion},
  whose notation it reuses. Every honesty flag of the notes is carried
  forward: sketches are labeled, quoted-standard facts are flagged,
  the status ledger is reproduced in Section~7, and no claim below
  exceeds what the underlying notes established.}

\begin{abstract}
For the pure cubic Thue equation $x^3 - d\,y^3 = m$, the companion
paper gave a deterministic $2^{O(s)}$ decision algorithm. Here we ask
what \emph{certifies} solvability. The certificate collapses to a
single object: a Thiel-format compact representation of the solution
element $\beta = x - y\alpha$ itself. Integrality and the norm
equation $N(\beta) = m$ verify deterministically in polynomial time
--- the latter by multiplicativity of norms on power products, the
mechanism behind Lagarias's Pell certificates and Ge's equality tests.
Membership in \textbf{MA} follows unconditionally. Membership in
\textbf{NP} is equivalent, modulo this architecture, to one newly
isolated problem --- the \textbf{Compact Vanishing Problem} (CVP):
decide deterministically whether a fixed coordinate functional
vanishes at a compactly represented algebraic integer. We classify
CVP into the sum-of-square-roots family, prove that ideal-theoretic
and radical-rigidity methods cannot express it, and explain by a
vacuity remark why Lagarias's degree-2 certificates never met it.
\end{abstract}

\maketitle

%% ==================================================================
\section{Introduction}\label{sec:intro}
%% ==================================================================

\subsection{From decision to certification}\label{sec:fromdecision}

The companion paper \cite{companion} proved that pure cubic Thue
equations $x^3 - d\,y^3 = m$ ($d \ge 2$ cube-free, not a cube;
$m \ne 0$; input size $s = O(\log d + \log|m|)$ in binary) are
decided, with full solution lists, by a deterministic unconditional
algorithm in time $2^{O(s)}$ --- a first worst-case bound of this
shape for a Thue stratum of Smale's fifth problem \cite{Sma98}. Its
Remark~7.5 named the three obstructions to \emph{polynomial-time}
decision: the unit (the principal ideal problem, classically believed
hard, polynomial only quantumly \cite{Hal07}); the zero test
(certified zero-testing against $2^{O(s)}$-digit height bounds); and
the window (quasi-linear in $R = 2^{\Theta(s)}$
(\cite[Lemma~C]{companion}, after Sha \cite{Sha19}) unless the strong
window conjecture is proven).

The natural target below polynomial time is therefore \textbf{NP
membership}: poly-size certificates, deterministic poly-time
verification. This paper is that pivot, and its outcome is
asymmetric: two thirds of the verifier's work is classical and easy,
and the entire remaining difficulty isolates into one new problem.

\subsection{One rung down: Lagarias's Pell certificates}
\label{sec:lagarias}

Lagarias \cite{Lag79} proved that solvability of binary quadratic
Diophantine equations over $\mathbb{Z}$ is in NP, although the
fundamental solutions have exponentially many digits: the certificate
is a power-product representation, and both verifiable conditions ---
integrality and the norm equation --- are \emph{multiplicative},
hence compatible with power products. (Negative Pell decision is even
in NP $\cap$ coNP by the same technology, with no classical
polynomial algorithm known.) This paper asks what happens to that
magic one rung up, at degree 3.

\subsection{The two main results, and the crux they isolate}
\label{sec:results}

The certificate architecture (Section~\ref{sec:cert}) simplifies an
outside reviewer's proposed three-part scheme (unit + reduced
representative + orbit index): all three collapse into the single
object $\beta$, the solution element itself in compact
representation. Verification splits into three steps --- integrality,
norm, Thue shape --- and:

\begin{itemize}
\item \textbf{Integrality and norm verify deterministically in
  polynomial time} (Section~\ref{sec:cert}), by the format's
  denominator bookkeeping and by norm multiplicativity on power
  products, the latter a special case of Ge's deterministic
  multiplicative-relation testing \cite{Ge93}.
\item \textbf{Proposition~\ref{prop:ma} (MA membership,
  unconditional).} Pure cubic Thue solvability is in MA, with
  one-sided error: the remaining condition --- the vanishing of one
  coordinate --- is testable modulo random primes, with the prime
  range calibrated against an explicit height bound extracted from
  the compact data.
\item \textbf{The crux (Section~\ref{sec:cvp}).} The single
  non-derandomized step is named: \textbf{CVP, the Compact Vanishing
  Problem} --- decide deterministically whether a fixed
  $\mathbb{Q}$-linear coordinate functional vanishes at a compactly
  represented algebraic integer. Cubic Thue $\in$ NP
  \textbf{iff-modulo} CVP (Corollary~\ref{cor:npmodcvp}, with the
  honest gloss given there); conditionally, under derandomization
  hypotheses implying MA = NP, membership in NP holds today.
\item \textbf{The vacuity remark (Remark~\ref{rem:vacuity}).} At
  degree 2 the module $\mathbb{Z} + \mathbb{Z}\sqrt{d}$ is \emph{all}
  of $\mathbb{Z}[\sqrt{d}]$: the shape condition is vacuous and CVP
  never arises. The coordinate condition is genuinely new at degree 3
  --- CVP, not the unit computation, is the honest frontier of the NP
  question.
\end{itemize}

Section~\ref{sec:toolbox} explains why CVP resists the classical
toolbox: four attack routes --- Thiel/Ge equality testing,
Bl\"omer's radical sums, ideal-theoretic reformulations, the Baker
route --- are each shown inapplicable, for one unified reason.

\subsection{Organization}\label{sec:org}

Section~\ref{sec:cert} fixes the certificate and verifier;
Section~\ref{sec:ma} proves MA membership; Section~\ref{sec:cvp}
states CVP, its three faces, and the vacuity remark;
Section~\ref{sec:toolbox} is the impossibility ledger;
Sections \ref{sec:fragments}--\ref{sec:open} collect positive
fragments, open problems, and the status ledger.

%% ==================================================================
\section{The certificate and the verifier}\label{sec:cert}
%% ==================================================================

\subsection{Notation (from the companion)}\label{sec:notation}

$d = ab^2 \ge 2$ cube-free, not a cube, $a, b$ coprime squarefree;
$\alpha = \sqrt[3]{d}$; $K = \mathbb{Q}(\alpha)$, complex cubic, one
real embedding $\sigma_1$ and a conjugate pair $\sigma_2, \sigma_3$;
unit group $\langle -1, \varepsilon \rangle$ of rank one, regulator
$R = \log \sigma_1(\varepsilon) = 2^{O(s)}$
\cite[\S~2.3]{companion}. For $g \in K$ write
$g = c_0(g) + c_1(g)\alpha + c_2(g)\alpha^2$; the index
$[O_K : \mathbb{Z}[\alpha]]$ divides $3b$, so $3b\,c_i(g) \in
\mathbb{Z}$ for $g \in O_K$ \cite[\S~2.2]{companion}. A solution
$(x, y)$ of $x^3 - dy^3 = m$ is the element $\beta = x - y\alpha \in
\mathbb{Z}[\alpha]$ with $N(\beta) = m$ and \textbf{Thue shape}
$z(\beta) := c_2(\beta) = 0$. We write $M := \mathbb{Z} +
\mathbb{Z}\alpha$ for the rank-2 module in which solutions live.

\subsection{Compact representations}\label{sec:compact}

A \textbf{compact representation} (Thiel format) of $\beta \in O_K$
is a power product
\[
\beta \;=\; \beta_0 \cdot \prod_{i=1}^{T} \beta_i^{2^i},
\]
with $T = O(\log h(\beta))$ factors, each $\beta_i \in K$ given by
$\mathrm{poly}(s)$ bits of coordinate and denominator data. The
existence and size bounds are Thiel's, \textbf{unconditional and in
arbitrary degree} \cite[Def.~6.3.1, Cor.~6.3.9--6.3.10]{Thi95}: every
unit and every generator of a principal ideal admits a compact
representation with $T \le \log(\log n + (n-1)\log H) + 2 =
O(\log h(\beta))$ factors of $\mathrm{poly}(\log|\Delta|)$ bits each.
(Thiel's normal form uses descending exponents $2^{T-i}$ and rational
factors $\alpha_i/d_i$, $d_i \le \sqrt{\Delta}$ --- cosmetically, not
mathematically, different from the form displayed above; the leading
factor additionally carries $O(\log|N(\beta)|)$ bits, harmless here.)
Efficient computation from power-product data and log-embedding
approximations is Biasse--Fieker \cite[Alg.~7, Prop.~5.1]{BF14},
unconditional given its input. The quadratic antecedent is
\cite{BTW95}.

For our instances a solution element has $\log$-height
$O(R + \log|m|) = 2^{O(s)}$ (unit reduction and the two-sided window
of \cite[Lemmas U and C]{companion}), so $T = O(s)$: the certificate
has $\mathrm{poly}(s)$ size. Existence is \emph{constructive}: the
honest prover runs the companion's exponential algorithm once; the
tracked-generator walk of \cite[Lemma~A$'$]{companion}, compacted
into Thiel's normal form via \cite[Alg.~7]{BF14}, emits exactly this
format. (Schoof's Arakelov machinery \cite{Sch08} supplies the
reduced-divisor arithmetic used elsewhere but contains no
compact-representation theorem --- a citation-scope point verified
against the source.)

\subsection{The certificate is the solution element itself}
\label{sec:certisbeta}

\begin{quote}
\textbf{Certificate (YES instances).} A Thiel-format compact
representation of $\beta = x - y\alpha$.
\end{quote}

No $\varepsilon$, no reduced representative, no orbit index $k$, no
window: if $\beta$ is a solution, its compact representation alone
witnesses it. In particular the strong window conjecture of
\cite[Remark~7.4]{companion} is \emph{not needed} for
YES-certificates at all; it re-enters only for the prover-efficiency
narrative and on the coNP side (Section~\ref{sec:windowconj}).

\subsection{The verifier, step by step}\label{sec:verifier}

\begin{enumerate}
\item \textbf{V1 --- Integrality}: $\beta \in O_K$, and $\beta \in
  \mathbb{Z}[\alpha]$ after the $3b$-bookkeeping. Deterministic
  polynomial time: the format carries denominator data checked by
  exact ideal arithmetic on $\mathrm{poly}$-size objects; and since
  coordinate denominators in $O_K$ divide $3b$, membership in
  $\mathbb{Z}[\alpha]$ is a \emph{congruence} condition modulo $3b$,
  computable by repeated squaring in a $\mathrm{poly}(s)$-size
  quotient ring. \emph{(Mechanism sketch; the format-level details
  are Thiel's.)}
\item \textbf{V2 --- Norm}: $N(\beta) = m$. Deterministic polynomial
  time, because \textbf{norms are multiplicative on power products}:
  \[
  N(\beta) \;=\; N(\beta_0)\cdot\prod_{i=1}^{T}
  N(\beta_i)^{2^i},
  \]
  each $N(\beta_i)$ a small explicit rational, exponents in binary
  --- exactly the magic that powered Lagarias's Pell certificates
  \cite{Lag79}. Verifying the exponent identity without factoring is
  standard multiplicative technology: coprime-basis bookkeeping
  \cite{vG25}, or a special case of Ge's multiplicative-relation test
  \cite{Ge93}; the sign is read off certified embedding data.
  \emph{(Sketch; the point is that every ingredient is
  multiplicative.)}
\item \textbf{V3 --- Thue shape}: $z(\beta) = 0$. \textbf{This is the
  entire remaining difficulty.} It is testable \emph{randomly} in
  polynomial time (Section~\ref{sec:ma}); whether
  \emph{deterministically} is precisely CVP (Section~\ref{sec:cvp}).
\end{enumerate}

Completeness of the architecture is immediate: $\beta \in
\mathbb{Z}[\alpha]$ with $N(\beta) = m$ and $c_2(\beta) = 0$ \emph{is}
a solution $(x, y) = (c_0(\beta), -c_1(\beta))$, and conversely.

%% ==================================================================
\section{Merlin--Arthur membership}\label{sec:ma}
%% ==================================================================

\begin{proposition}[MA membership, unconditional]\label{prop:ma}
Pure cubic Thue solvability is in MA, with perfect completeness and
one-sided, amplifiable soundness error.
\end{proposition}

\begin{longproof}{Proof (complete modulo the compact-representation
format bounds quoted in Section~\textup{\ref{sec:compact}})}
Merlin sends the compact representation of a claimed solution element
$\beta$. Arthur runs V1 and V2 of Section~\ref{sec:verifier}
deterministically, then tests V3 modulo random primes.

\textbf{The integer to test.} Since $\alpha, \beta \in O_K$, the
trace $W := \mathrm{Tr}_{K/\mathbb{Q}}(\beta\alpha)$ is a rational
integer, and by the trace identity of Section~\ref{sec:traceface},
\[
z(\beta) = 0 \iff W = 0.
\]

\textbf{Height bound from the compact data.} Embeddings are
multiplicative on the format: for each $j$,
\[
\log|\sigma_j(\beta)| \;=\; \log|\sigma_j(\beta_0)| +
\sum_{i=1}^{T} 2^i \log|\sigma_j(\beta_i)|.
\]
Each factor carries $\mathrm{poly}(s)$ bits, so
$\log|\sigma_j(\beta_i)| \le P(s)$ for an explicit polynomial $P$,
and $T = O(s)$ gives $\log_2\mathrm{house}(\beta) \le P(s)\,2^{T+1}
\le 2^{c_3 s}$ for an explicit $c_3$. Hence
$|W| \le 3\,\mathrm{house}(\beta)\,\mathrm{house}(\alpha)$, so
$\log_2 |W| \le 2^{c_3 s} + O(s) \le 2^{c_4 s}$; if $W \ne 0$, its
number of distinct prime divisors obeys $\omega(W) \le \log_2|W| \le
2^{c_4 s}$.

\textbf{The prime range for a $2/3$ threshold.} Let $E$ be the set of
primes dividing $3d$ or any denominator in the format;
$\#E \le \mathrm{poly}(s)$. Arthur draws $p$ uniformly from the
primes in $[17, 2^t]$: by Chebyshev-type bounds ($\pi(x) > x/\ln x$
for $x \ge 17$), the choice $t = c_4\,s + O(\log s)$, with explicit
absolute constants, supplies at least $3\,(\omega(W) + \#E)$ primes,
so $\Pr[\,p \in E \text{ or } p \mid W\,] \le 1/3$. If $p \in E$,
Arthur accepts vacuously (inside the $1/3$). Otherwise he computes
$W \bmod p$ in polynomial time: writing multiplication-by-$\beta_i$
as $3\times 3$ rational matrices $M_{\beta_i}$ in the power basis, he
evaluates $M_\beta \equiv M_{\beta_0}\prod_i M_{\beta_i}^{2^i}
\pmod p$ by repeated squaring, then
$W \equiv \mathrm{tr}(M_\beta M_\alpha) \bmod p$, accepting iff
$W \equiv 0$.

\textbf{Completeness.} For a YES instance Merlin sends a true
solution element: V1, V2 pass deterministically and $W = 0$, so
$W \equiv 0$ modulo \emph{every} prime --- Arthur accepts with
probability 1.

\textbf{Soundness (one-sided).} For a NO instance, any certificate
passing V1 and V2 must have $z(\beta) \ne 0$ --- otherwise $\beta$
would be a solution --- hence $W \ne 0$, and Arthur rejects with
probability $\ge 2/3$. Arthur never rejects a valid certificate.

\textbf{Amplification.} $k$ independent primes give error $(1/3)^k$;
widening the range to $2^{s^c}$, $c \ge 2$, already makes a single
prime's error exponentially small, since $\omega(W)/\pi(2^{s^c}) \le
2^{c_4 s - s^c}\,\mathrm{poly}(s)$.
\end{longproof}

\emph{Complexity-class bookkeeping.} Step V3 alone is a coRP test:
vanishing instances are accepted with probability 1, non-vanishing
instances rejected with high probability. The outer existential
quantifier over certificates therefore places the decision problem in
$\exists\cdot\mathrm{coRP} = \mathrm{MA}$ --- which is the precise
sense of Proposition~\ref{prop:ma}.

\emph{(Provenance: the notes' ledger listed this as
``proposition-with-sketch, needs a careful write-up (height bound on
$z(\beta)$ from the compact data; prime-range constants)''. The
bounds above discharge those two items; the residual dependence is
the quoted format bounds of Section~\ref{sec:compact}.)}

%% ==================================================================
\section{The Compact Vanishing Problem}\label{sec:cvp}
%% ==================================================================

\subsection{Statement and consequence}\label{sec:cvpstatement}

\begin{quote}
\textbf{$\mathrm{CVP}_K$ (Compact Vanishing Problem).} Given a
compact representation of an algebraic integer $\beta$ in a fixed
cubic field $K$, decide in deterministic polynomial time whether a
fixed $\mathbb{Q}$-linear coordinate functional vanishes at $\beta$.
\end{quote}

For our instances the functional is $c_2$, the $\alpha^2$-coordinate
in $K = \mathbb{Q}(\sqrt[3]{d})$. (The acronym clash with the lattice
Closest Vector Problem is noted; no lattice problem appears here.)

\begin{corollary}[NP iff-modulo CVP]\label{cor:npmodcvp}
If $\mathrm{CVP}_K \in$ \textup{P} for the instances arising from
Thue certificates, then pure cubic Thue solvability is in
\textup{NP}, unconditionally. Conversely --- the ``iff-modulo''
gloss, architectural rather than formal --- CVP is the \emph{only}
step of the canonical verifier that is not deterministic polynomial
time: any \textup{NP} proof through this certificate collapses to
deciding CVP on these instances. Under standard derandomization
hypotheses (which give \textup{MA} $=$ \textup{NP}), \textup{NP}
membership holds conditionally today.
\end{corollary}

The notes established three equivalent faces of CVP: a trace form, a
lattice-membership form, a bounded-width branching-program form.

\subsection{Face one: the trace form}\label{sec:traceface}

For $\beta = x + y\alpha + z\alpha^2$ one computes $\beta\alpha =
zd + x\alpha + y\alpha^2$, and $\mathrm{Tr}(1) = 3$,
$\mathrm{Tr}(\alpha) = \mathrm{Tr}(\alpha^2) = 0$ give
\[
\mathrm{Tr}(\beta\alpha) \;=\; 3d\,z, \qquad\text{i.e.}\qquad
z(\beta) \;=\; \frac{\mathrm{Tr}(\beta\alpha)}{3d}.
\]
CVP for our functional is exactly: \emph{is $\mathrm{Tr}(\beta\alpha)
= 0$ for compact $\beta$?} The coordinate functionals \emph{are}
trace forms --- the dual basis of the power basis under the trace
pairing gives $c_2(\beta) = \mathrm{Tr}(w\beta)$ with
$w = \alpha/(3d)$, and likewise for $c_1, c_0$ (identities
machine-checked during this draft's assembly).

\subsection{Face two: plane membership for Arakelov-pinned elements}
\label{sec:planeface}

The second face rests on a small clean result, the micro-yield of the
notes' Ge-lattice attempt.

\setcounter{lemma}{2}% match the notes' numbering: Lemma 4.3
\begin{lemma}[Arakelov pinning]\label{lem:pinning}
From a valid compact representation of $\beta \in O_K$, a
deterministic polynomial-time verifier computes the full
\textbf{Arakelov divisor} of $\beta$: the exact ideal $I = (\beta)$
in Hermite normal form, and the log-embeddings
$\log|\sigma_j(\beta)|$ to any requested $\mathrm{poly}(s)$
precision. This data determines $\beta$ up to sign.
\end{lemma}

\begin{proof}[Proof sketch]
The ideal: clear the format's denominators first and work with the
numerator ideal of the compacted element --- so that $I$ is integral,
as the HNF intersection $L = I \cap M$ of
Section~\ref{sec:planeface} strictly requires --- then reduce the
power product by repeated squaring on reduced ideals with tracked
distances --- infrastructure arithmetic is precisely this machinery
\cite{BW88,Sch08} --- every intermediate object staying
$\mathrm{poly}$-size. The logs: interval arithmetic on the factors'
certified embedding data, via the exact telescoping of
Section~\ref{sec:ma}. Uniqueness: if $(\beta) = (\beta')$ with equal
log vectors, $\beta/\beta'$ is a unit with vanishing log-embeddings,
hence torsion, hence $\pm 1$ (the torsion of $K$ is $\{\pm 1\}$ by
the real embedding). \emph{(sketch)}
\end{proof}

So the certificate pins a \emph{unique} candidate (up to sign,
harmless since $z(-\beta) = -z(\beta)$): \textbf{compact certificates
verify everything about $\beta$ except one plane membership.} Indeed
$z(\beta) = 0 \iff \beta \in M$, and $L := I \cap M$ is an explicit
rank-2 lattice with a $\mathrm{poly}$-size basis (HNF intersection).
Hence:

\begin{quote}
\textbf{CVP (Thue instances) $\iff$ does the element pinned by a
given Arakelov divisor lie in an explicit rank-2 sublattice?}
\end{quote}

\subsection{Face three: width-3 ABP trace zero-testing}
\label{sec:abpface}

Writing multi\-pli\-ca\-tion-by-$\beta_i$ as $3\times 3$ integer
matrices (after clearing the format's denominators), the compact
representation becomes a \textbf{matrix word with repeated-squaring
structure}, and by the trace form CVP reads: \emph{decide whether the
trace of a width-3 integer matrix word is zero}. That is polynomial
identity testing for width-3 algebraic branching programs --- a
highly structured, whitebox, repeated-squaring instance, and the
right literature door: it connects the NP question for cubic Thue
equations to bounded-width PIT derandomization by the shortest path
we know of (Section~\ref{sec:equslp}). One structural advantage must
be stated explicitly: the matrices $M_{\beta_i}$ all lie in the
commutative subalgebra $\mathbb{Q}[M_\alpha] \cong K$ of $3\times 3$
matrices --- the word is not a generic non-commutative width-3 ABP
but one over a \textbf{commutative} matrix ring, a class
significantly better understood in the PIT literature and
correspondingly a more tempting target for derandomization.

\subsection{Remark (Lagarias vacuity: why 1979 stopped at degree 2)}
\label{sec:vacuity}

\setcounter{remark}{4}% match the notes' numbering: Remark 4.5
\begin{remark}\label{rem:vacuity}
In the quadratic case the module $\mathbb{Z} + \mathbb{Z}\sqrt{d}$ is
\emph{all} of $\mathbb{Z}[\sqrt{d}]$: the Thue-shape condition is
vacuous, CVP never arises, and norm-checking alone suffices --- which
is why Lagarias's certificates \cite{Lag79} close at degree 2 with
purely multiplicative technology. The coordinate condition is
genuinely \textbf{new at degree 3}: CVP is exactly what separates the
cubic certificate problem from Pell --- this, not the unit
computation, is the honest frontier of the NP question, and the
explanation of the shift from degree 2 (NP since 1979) to degree 3
(MA now).
\end{remark}

%% ==================================================================
\section{Why the classical toolbox stops}\label{sec:toolbox}
%% ==================================================================

We place CVP on the map, then close four attack routes --- distilling
the notes' 1991--2026 literature sweep and recorded attempts. No
route is closed by a hardness theorem; each is closed by structural
inapplicability, a different and more instructive thing.

\setcounter{subsection}{-1}% the notes number this part 5.0--5.5

\subsection{Where CVP sits: the sum-of-square-roots family}
\label{sec:family}

The cleared trace value admits a $\mathrm{poly}(s)$-size
straight-line program (repeated squaring), so CVP is an instance of
EquSLP --- zero-testing SLP-given integers --- in coRP (the
random-prime test of Section~\ref{sec:ma}; the observation goes back
to Sch\"onhage \cite{Sch79}), with its positivity cousin PosSLP in
the counting hierarchy \cite{ABKM09,ESY14}. Deterministic polynomial
time is open for the whole family, whose totem is the
sum-of-square-roots problem: poly-time for explicitly given radical
sums \cite{Blo91}, uniform TC$^0$ for deciding whether a rational
radical sum is rational \cite{HBMOW10}, the separation-bound frontier
only recently moved by the subspace theorem \cite{EHS24}. The notes'
verdict, adopted here: CVP is \textbf{genuinely new, of
sum-of-square-roots-class flavor} --- in coRP, in CH, deterministic
case open --- with one unexplored asset recorded in
Section~\ref{sec:bridge}.

\subsection{The additive/multiplicative divide}\label{sec:divide}

Thiel-style deterministic \emph{equality} testing of compact
representations is a \textbf{multiplicative} technology:
$\beta = \gamma$ reduces to the quotient being $1$, and units
$\ne \pm1$ are separated from $1$ by Dobrowolski-type height lower
bounds \cite{Dob79} on the log-embedding lattice --- polynomial
precision suffices. Ge's theorem \cite{Ge93} (see
\cite[Thm.~165]{vG25}) extends this: \emph{multiplicative relations}
among power products are decidable in deterministic polynomial time
--- exactly why the norm step V2 is easy. The compact-representation
toolbox stops at the trace, precisely: \cite{BTW95} adds two-term
sign comparison in real quadratic fields, and \textbf{no additive
$\ge 3$-term test exists anywhere in the literature}.

CVP is \textbf{additive}: $\mathrm{Tr}(\beta\alpha)$ is a sum of
three conjugate terms. Separation is not the problem --- integrality
gives it for free (the cleared value is a rational integer, so
nonzero means $\ge 1$). \textbf{Evaluation is the problem}:
certifying the sum numerically needs absolute precision comparable to
the archimedean spread
$\max_j \log|\sigma_j(\beta)| - \min_j \log|\sigma_j(\beta)|$, and
for \emph{true} Thue solutions the spread is $2^{\Theta(s)}$ forced:
a solution has $|\sigma_1(\beta)| = |m|/|\sigma_2(\beta)|^2$ tiny
\emph{precisely because} it is a solution. The easy
``balanced-$\beta$'' case of CVP can therefore never contain the
instances we care about. The angle formulation hits the same wall:
verifying $\arg \sigma_2(\beta\alpha) \approx \pm\pi/2$ to error
$e^{-\mathrm{spread}}$ needs the small factors' arguments to
exponentially many bits.

\textbf{The conjugate-certificate attempt, and its circle.} One can
move to the Galois closure $L = K(\omega)$ and note
$\mathrm{Tr}(\beta\alpha) = 0 \iff \beta + \omega\sigma_2(\beta) +
\omega^2\sigma_3(\beta) = 0$; certifying compact representations of
the conjugates is fine, but the final check is again an additive
vanishing of compact objects --- the divide reappears intact.
\textbf{Sums of compact representations are not compact; that
sentence is the obstacle.}

\subsection{Bl\"omer is structurally inapplicable, not merely
insufficient}\label{sec:blomer}

Bl\"omer's deterministic polynomial-time algorithms for sums of
radicals \cite{Blo91,Blo93} might look like the natural import: our
elements are $\mathbb{Q}$-combinations of $1, d^{1/3}, d^{2/3}$. But
his determinism rides on \emph{Siegel-type rigidity}: independent
radicals admit \textbf{no} nontrivial $\mathbb{Q}$-linear relations,
so zero-testing reduces to coefficient collection on explicit
poly-size coefficients. Our trace-zero instances \textbf{are}
nontrivial relations in the rank-2 trace-zero module --- the exact
phenomenon his theorem excludes --- with coefficients $x, y, z$
compact rather than explicit. Both pillars --- rigidity and explicit
arithmetic --- fail simultaneously.

\subsection{No ideal-theoretic route: the hyperplane argument}
\label{sec:hyperplane}

The tempting program ``tear into Thiel's ideal equality tests to
break CVP'' is closed before it opens, by rank counting.

\begin{proposition}\label{prop:hyperplane}
No ideal containment, equality, or divisibility statement in $O_K$
can express the condition $z(\beta) = 0$, even in principle.
\end{proposition}

\begin{proof}
The trace-zero condition $\mathrm{Tr}(\beta\alpha) = 0$ defines a
$\mathbb{Q}$-\textbf{hyperplane} in $K$ (the coordinate functionals
are trace forms, Section~\ref{sec:traceface}). Ideals are
\textbf{full-rank} $\mathbb{Z}$-modules: any nonzero ideal of $O_K$
has rank 3, while any module contained in the hyperplane has rank
$\le 2$.
\end{proof}

Thiel/Ge machinery tests relations among full-rank multiplicative
objects; CVP lives in a corank-1 additive slice invisible to all of
them. The bridge of Section~\ref{sec:wall} is not hiding inside ideal
arithmetic; it must be genuinely new.

\subsection{The EquSLP caution, and the Baker route's ceiling}
\label{sec:equslp}

Two doors from complexity theory, each honestly measured:

\begin{itemize}
\item \textbf{Iterated squaring is the hard core.} ABKM
  \cite[Prop.~2.1]{ABKM09} show that EquSLP's difficulty concentrates
  on repeated-squaring instances --- our shape --- and derandomizing
  such identity tests implies circuit lower bounds \cite{KI04}. This
  is \textbf{a caution, not a hardness proof}, for our width-3 single
  words: CVP could be easy without lower-bound consequences, but the
  precedent says the general-toolbox route is priced.
\item \textbf{The Baker route has a height ceiling.}
  Galby--Ouaknine--Worrell \cite{GOW15} decide sign and zero
  questions for entries of $M^n$, $\dim \le 3$, in polynomial time
  --- for \textbf{poly-height} inputs: the engine is lower bounds for
  linear forms in logarithms \cite{Mat00}, whose constants stay
  polynomial only while heights and the number of logarithms stay
  bounded. Thue certificates have
  $h(\varepsilon) = R/3 = 2^{\Theta(s)}$; at $\Theta(s)$ logarithms
  the Baker constants degrade exponentially --- the same wall that
  pins LRS Positivity at order $\ge 6$ \cite{OW14}. The bounded
  fragment this route \emph{does} give is
  Proposition~\ref{prop:bounded}.
\end{itemize}

\subsection{One wall, four costumes; the missing tool, named}
\label{sec:wall}

The four failures above are one wall:

\begin{enumerate}
\item \textbf{Trace non-multiplicativity}
  (Section~\ref{sec:divide}).
\item \textbf{Sums of compacts are not compact}
  (Section~\ref{sec:divide}).
\item \textbf{The archimedean spread} (Section~\ref{sec:divide}):
  $2^{\Theta(s)}$ precision forced exactly on true solutions.
\item \textbf{Plane non-stability} (the notes' Ge-lattice attempt):
  the rescaling trick --- bring the thin strip
  $\{|x - y\alpha| \sim e^{\ell_1}\}$ to bounded size by
  $\varepsilon^{-j}$ --- dies because $M$ is not stable under units:
  $\varepsilon^{-j}M \ne M$. Lattice-point location in the plane is
  polynomial \emph{at polynomial scales}, but our scales are
  $e^{2^{O(s)}}$, described only by their logs.
\end{enumerate}

What a solution needs, stated once: a
\textbf{multiplicative-to-additive bridge} --- any theorem letting
membership in a fixed rank-2 plane be tested through data carried
multiplicatively (relation lattices, Arakelov classes, unit orbits).
Nothing in the 1991--2026 literature provides one; building it is the
actual open problem, and by Section~\ref{sec:planeface} it now has a
precise geometric formulation.

%% ==================================================================
\section{Positive fragments}\label{sec:fragments}
%% ==================================================================

\subsection{The bounded fragment: CVP at bounded rank and small
height}\label{sec:bounded}

\begin{proposition}[bounded CVP --- stated with a proof plan]
\label{prop:bounded}
Restrict $\mathrm{CVP}_K$ to instances whose factor list
$\{\beta_0, \ldots, \beta_T\}$ generates a subgroup of
$K^\times/\{\pm1\}$ of multiplicative rank $O(1)$ admitting
generators of Weil height $\mathrm{poly}(s)$. This fragment is
decidable in deterministic polynomial time.
\end{proposition}

\begin{proof}[Proof plan (honestly flagged: a plan, not a proof)]
(1) Compute the relation lattice of the factors --- deterministic
polynomial time by Ge \cite{Ge93}, \cite[Thm.~165]{vG25} --- and
rewrite the word over $O(1)$ independent poly-height generators with
binary exponents. (2) The trace becomes a fixed short sum of
conjugate power products in $O(1)$ poly-height bases. (3) Zero-test
via the Baker--Matveev machinery as in the dimension-3 matrix-power
analysis of \cite{GOW15}: with $O(1)$ logarithms of poly-height
numbers and binary exponents, Matveev's bounds \cite{Mat00} are
polynomial --- noting that Matveev's constant degrades exponentially
in the number $k$ of logarithms, which is exactly why the $O(1)$
multiplicative-rank restriction is load-bearing and cannot be
relaxed, so certified interval evaluation at $\mathrm{poly}(s)$
precision decides against the free integrality separation. Steps
(2)--(3) need the same case-discipline GOW exercise at dimension 3
and are not written out here.
\end{proof}

\emph{What this fragment is and is not.} It is the notes' verdict
``publishable-but-bounded''. It does \textbf{not} cover Thue
certificates: their factor lists contain $\varepsilon$ (or its walk
factors), of height $R/3 = 2^{\Theta(s)}$ --- precisely the ceiling
of Section~\ref{sec:equslp}. No claim of progress on the Thue NP
question is made here.

\subsection{The strong window conjecture: prover efficiency and coNP}
\label{sec:windowconj}

The strong window conjecture \cite[Remark~7.4]{companion} --- every
Thue-shape index of a unit-reduced representative satisfies
$|k| \le C(1 + \log|m|/R)$, hence $O(s)$ by the regulator floor
\cite{ADF16} --- turned out to be unnecessary for YES-certificates
(Section~\ref{sec:certisbeta}). Its remaining roles: \textbf{prover
efficiency} --- the honest prover runs the companion's algorithm once
($2^{O(s)}$); the conjecture caps the final scan at $O(s)$ indices
per representative, though the infrastructure walk remains
$\widetilde{\Theta}(R)$, so it trims the narrative, not the
exponential --- and \textbf{the coNP side}, where certifying a NO
instance would require certifying every ideal class's orbit window
zero-free, and the conjecture caps the indices to certify per orbit
at $\mathrm{poly}(s)$, making class-enumeration NO-certificates
conceivable. \emph{(Ledger flag, preserved: the coNP side is
untouched in this work.)}

%% ==================================================================
\section{Open problems}\label{sec:open}
%% ==================================================================

\subsection{CVP itself}\label{sec:cvpopen}

Is $\mathrm{CVP}_K$ --- even just for Thue-shaped instances --- in P?
The notes' two probed routes are blocked, not dead: the \emph{$p$-adic
route} ($\mathrm{ord}_p$ of power products is poly-computable per
prime) stumbles on the fact that valuations do not see addition; the
\emph{canonical-form route} (Thiel equality of $\beta$ against its
projection) needs a compact representation of the projection, and
sums of compacts are not compact. The sharpest open formulation is
the ABP face: \textbf{is the repeated-squaring width-3 matrix-word
class among the deterministically testable bounded-width PIT cases?}
A positive answer completes: cubic Thue $\in$ NP, unconditionally ---
Lagarias one rung up, for real.

\subsection{The multiplicative-to-additive bridge}\label{sec:bridge}

Build any theorem letting membership in a fixed rank-2 plane be
tested through multiplicatively carried data
(Section~\ref{sec:wall}). The one unexplored asset, recorded in the
notes' literature verdict: our three summands are Galois-conjugate,
and the bases' full multiplicative relation lattice is
\emph{computable} (Ge). The additive analog of
\cite[Lemma~166]{vG25} --- detecting trace-zero from the relation
lattice plus unit-lattice geometry --- has never been attempted. That
is the door.

\subsection{DIVIS hardness}\label{sec:divis}

The program's synthesis records that the divisibility language
$\mathrm{DIVIS} = \{(A, B) \in \mathbb{Z}[x]^2 : \exists x \in
\mathbb{Z},\ A(x) \mid B(x)\}$ is in NP (two proofs, one the
expedition's, by an elementary integer-remainder argument --- Runge
in miniature); its hardness is unstudied. It is the nearest candidate
for a hardness foothold adjacent to the Thue stratum.

\subsection{The coNP certificate question}\label{sec:conp}

Is pure cubic Thue solvability in coNP? Degree 2 says yes one rung
down (negative Pell is in NP $\cap$ coNP \cite{Lag79}). At degree 3
the ingredients on the table are the strong window conjecture and
class-enumeration certificates (Section~\ref{sec:windowconj});
nothing beyond the framing exists yet.

\subsection*{Status ledger (carried from the notes, unchanged in
spirit)}

\begin{itemize}
\item MA membership: proposition with the height-bound and
  prime-range write-up supplied (Section~\ref{sec:ma}); the
  format-bound dependence is now DISCHARGED by the pinned citations
  (Section~\ref{sec:compact}).
\item Compact-representation existence and size: \textbf{debt paid}
  --- Thiel's thesis read in full (Def.~6.3.1, Cor.~6.3.9--6.3.10,
  unconditional, arbitrary degree), efficient computation pinned to
  \cite[Alg.~7, Prop.~5.1]{BF14}; the erroneous compaction citation
  to \cite{Sch08} corrected (it contains no compact representations
  --- verified against the source). The remaining ledger flags may
  now be stripped at final submission.
\item CVP: open, named, isolated --- the right-sized next
  mathematics.
\item coNP side: untouched here.
\end{itemize}

%% ==================================================================
%% References
%% ==================================================================

\begin{thebibliography}{companion}

\bibitem[ABKM09]{ABKM09} E. Allender, P. B\"urgisser,
  J. Kjeldgaard-Pedersen, P. B. Miltersen, \emph{On the complexity of
  numerical analysis}, SIAM J. Comput. 38:5 (2009) 1987--2006.

\bibitem[ADF16]{ADF16} S. Astudillo, F. D\'iaz y D\'iaz, E. Friedman,
  \emph{Sharp lower bounds for regulators of small-degree number
  fields}, J. Number Theory 167 (2016) 232--258.

\bibitem[BF14]{BF14} J.-F. Biasse, C. Fieker, \emph{Subexponential
  class group and unit group computation in large degree number
  fields}, LMS J. Comput. Math. 17 (2014) 385--403. (\S~5,
  Algorithm~7, Proposition~5.1: efficient computation of compact
  representations; unconditional given its input.)

\bibitem[Blo91]{Blo91} J. Bl\"omer, \emph{Computing sums of radicals
  in polynomial time}, Proc. 32nd IEEE FOCS (1991) 670--677.

\bibitem[Blo93]{Blo93} J. Bl\"omer, \emph{Simplifying expressions
  involving radicals}, PhD thesis, Freie Universit\"at Berlin, 1993.

\bibitem[BTW95]{BTW95} J. Buchmann, C. Thiel, H. C. Williams,
  \emph{Short representation of quadratic integers}, in:
  Computational Algebra and Number Theory, Math. Appl. 325, Kluwer,
  1995, 159--185.

\bibitem[BW88]{BW88} J. Buchmann, H. C. Williams, \emph{On the
  infrastructure of the principal ideal class of an algebraic number
  field of unit rank one}, Math. Comp. 50 (1988) 569--579.

\bibitem[companion]{companion} The diophantine-fable expedition,
  \emph{Pure cubic Thue equations are solvable in deterministic
  exponential time}, companion manuscript,
  \texttt{papers/cubic-thue-exp.md} in this repository (draft of
  2026-07-22).

\bibitem[Dob79]{Dob79} E. Dobrowolski, \emph{On a question of Lehmer
  and the number of irreducible factors of a polynomial}, Acta
  Arith. 34 (1979) 391--401.

\bibitem[EHS24]{EHS24} F. Eisenbrand, M. Haeberle, N. Singer, \emph{An
  improved bound on sums of square roots via the subspace theorem},
  Proc. SoCG 2024, LIPIcs 293, art. 54.

\bibitem[ESY14]{ESY14} K. Etessami, A. Stewart, M. Yannakakis, \emph{A
  note on the complexity of comparing succinctly represented
  integers, with an application to maximum probability parsing}, ACM
  Trans. Comput. Theory 6:2 (2014), art. 9.

\bibitem[Ge93]{Ge93} G. Ge, \emph{Testing equalities of
  multiplicative representations in polynomial time}, Proc. 34th IEEE
  FOCS (1993) 422--426; and PhD thesis (\emph{Algorithms related to
  multiplicative representations of algebraic numbers}),
  U.C. Berkeley, 1993.

\bibitem[GOW15]{GOW15} E. Galby, J. Ouaknine, J. Worrell, \emph{On
  matrix powering in low dimensions}, Proc. STACS 2015, LIPIcs 30,
  329--340.

\bibitem[Hal07]{Hal07} S. Hallgren, \emph{Polynomial-time quantum
  algorithms for Pell's equation and the principal ideal problem},
  J. ACM 54:1 (2007), art. 4.

\bibitem[HBMOW10]{HBMOW10} P. Hunter, P. Bouyer, N. Markey,
  J. Ouaknine, J. Worrell, \emph{Computing rational radical sums in
  uniform TC$^0$}, Proc. FSTTCS 2010, LIPIcs 8, 308--316.

\bibitem[KI04]{KI04} V. Kabanets, R. Impagliazzo, \emph{Derandomizing
  polynomial identity tests means proving circuit lower bounds},
  Comput. Complexity 13 (2004) 1--46.

\bibitem[Lag79]{Lag79} J. C. Lagarias, \emph{Succinct certificates
  for the solvability of binary quadratic Diophantine equations},
  Proc. 20th IEEE FOCS (1979) 47--54; extended version
  arXiv:math/0611209.

\bibitem[Mat00]{Mat00} E. M. Matveev, \emph{An explicit lower bound
  for a homogeneous rational linear form in the logarithms of
  algebraic numbers. II}, Izv. Math. 64:6 (2000) 1217--1269.

\bibitem[OW14]{OW14} J. Ouaknine, J. Worrell, \emph{Positivity
  problems for low-order linear recurrence sequences}, Proc. SODA
  2014, 366--379.

\bibitem[Sch79]{Sch79} A. Sch\"onhage, \emph{On the power of random
  access machines}, Proc. ICALP 1979, LNCS 71, 520--529.

\bibitem[Sch08]{Sch08} R. Schoof, \emph{Computing Arakelov class
  groups}, in: Algorithmic Number Theory, MSRI Publ. 44, Cambridge
  University Press, 2008, 447--495.

\bibitem[Sha19]{Sha19} M. Sha, \emph{Effective results on the Skolem
  Problem for linear recurrence sequences}, J. Number Theory 197
  (2019) 228--249.

\bibitem[Sma98]{Sma98} S. Smale, \emph{Mathematical problems for the
  next century}, Math. Intelligencer 20:2 (1998) 7--15.

\bibitem[Thi95]{Thi95} C. Thiel, \emph{On the complexity of some
  problems in algorithmic algebraic number theory}, PhD thesis,
  Universit\"at des Saarlandes, Saarbr\"ucken, 1995 (Definition
  6.3.1, Corollaries 6.3.9--6.3.10; Theorem 6.4.1: principal-ideal
  certificates in NP, unconditional); journal companion: \emph{Short
  proofs using compact representations of algebraic integers},
  J. Complexity 11 (1995) 310--329; announcement: ANTS-I 1994, LNCS
  877, 234--247.

\bibitem[vG25]{vG25} D. M. H. van Gent (lectures by
  H. W. Lenstra), \emph{Polynomial-time algorithms in algebraic
  number theory}, arXiv:2502.19036. (Thm. 165 = Ge's theorem; Lemma
  166; coprime bases.)

\end{thebibliography}

\end{document}
